\section{Architecture}
\label{sec:architecture}

Let $X_\ell \in \mathbb{R}^{T \times d}$ be the token stream at layer $\ell$,
$E_\ell \in \mathbb{R}^{M_e \times d}$ the event-role register bank, and
$P_\ell \in \mathbb{R}^{M_p \times d}$ the primitive register bank. The normal
decoder stream is updated by grouped-query causal self-attention:

\begin{equation}
H_\ell = X_\ell + \Attn_{\mathrm{causal}}\!\left(\LN(X_\ell)\right).
\label{eq:token_update}
\end{equation}

The event-role stream reads from token states:

\begin{equation}
\widetilde{E}_\ell =
E_\ell +
\Attn_{E \leftarrow X}\!\left(
Q=\LN(E_\ell),\,
K=\LN(H_\ell),\,
V=\LN(H_\ell)
\right).
\label{eq:event_read}
\end{equation}

The token stream then receives an event-role signal $Z^E_\ell$:

\begin{equation}
Z^E_\ell =
\Attn_{X \leftarrow E}\!\left(
Q=\LN(H_\ell),\,
K=\LN(\widetilde{E}_\ell),\,
V=\LN(\widetilde{E}_\ell)
\right).
\label{eq:event_write}
\end{equation}

In the from-scratch architecture, the event injection can use token state,
register state, and their interaction:

\begin{equation}
H'_\ell = H_\ell + W_E \left[H_\ell \concat Z^E_\ell \concat
                    (H_\ell \odot Z^E_\ell)\right].
\label{eq:residual_injection}
\end{equation}

For warm-started pretrained decoders, this residual fusion is too aggressive:
the pretrained hidden state $H_\ell$ dominates the injection path. The validated
warm-start rule is register-only and detached:

\begin{equation}
H'_\ell = H_\ell + \alpha_E W_E \stopgrad(Z^E_\ell),
\quad \alpha_E \ll 1.
\label{eq:warmstart_injection}
\end{equation}

The primitive stream reads from the event-enriched token and event states:

\begin{equation}
\widetilde{P}_\ell =
P_\ell +
\Attn_{P \leftarrow X,E}\!\left(
Q=\LN(P_\ell),\,
K=\LN([H'_\ell; \widetilde{E}_\ell]),\,
V=\LN([H'_\ell; \widetilde{E}_\ell])
\right),
\label{eq:primitive_read}
\end{equation}

and writes primitive information back to tokens:

\begin{equation}
Z^P_\ell =
\Attn_{X \leftarrow P}\!\left(
Q=\LN(H'_\ell),\,
K=\LN(\widetilde{P}_\ell),\,
V=\LN(\widetilde{P}_\ell)
\right).
\label{eq:primitive_write}
\end{equation}

The primitive-enriched token state is

\begin{equation}
H''_\ell =
H'_\ell + W_P\left[H'_\ell \concat Z^P_\ell \concat
                  (H'_\ell \odot Z^P_\ell)\right],
\label{eq:primitive_injection}
\end{equation}

or, in warm-start-safe mode,

\begin{equation}
H''_\ell = H'_\ell + \alpha_P W_P \stopgrad(Z^P_\ell).
\label{eq:warmstart_primitive_injection}
\end{equation}

Finally, the layer output combines the base feed-forward network with
primitive-conditioned adapters:

\begin{equation}
X_{\ell+1} =
H''_\ell +
\operatorname{FFN}_{\mathrm{base}}(H''_\ell) +
\sum_{k=1}^{K} g_k(\widetilde{P}_\ell)\,
\operatorname{Adapter}_k(H''_\ell).
\label{eq:ffn_adapters}
\end{equation}

\subsection{Parameterization and Practical Cost}
\label{subsec:parameter_cost}

The warm-started configuration uses hidden size $d=1024$, 28 decoder layers,
context length 1024, 4 event registers, 6 argument registers, and 14 primitive
registers. Cross-stream attention is applied every four layers. The model uses
17 proto-role properties, six primitive-conditioned FFN adapters of rank 32, and
rank-32 vocabulary-pressure projections. The from-scratch 450M comparison
configuration uses the same hidden size with 16 layers and 457.9M parameters.
The Qwen3-compatible warm-start model loads 595.8M pretrained backbone
parameters; the remaining PAT-ER parameters include the extended vocabulary,
side-state streams, auxiliary heads, pressure heads, and interface components.

The register counts are small relative to the token context. At a cross-stream
layer, token-register attention adds a term proportional to
$O(T(M_e+M_a+M_p)d)$, while causal self-attention remains proportional to
$O(T^2d)$ for sequence length $T$. For $T=1024$ and 24 total registers, the
register-attention term is therefore not the dominant asymptotic cost. The
reported experiments measure model metrics and memory-feasible training gates;
they do not claim a hardware-independent wall-clock benchmark.

\subsection{Primitive Semantic Heads}
\label{subsec:semantic_heads}

We use \emph{semantic head} operationally: a fixed side-state channel and
readout trained against a characteristic semantic target. This is not a claim
that every individual register has been mechanistically interpreted. Each
primitive head is associated with a characteristic question:

\begin{equation}
\begin{array}{lll}
\text{modus ponens} &:& \text{from which verified antecedent does this follow?}\\
\text{syllogism} &:& \text{which chained rule licenses this conclusion?}\\
\text{contradiction} &:& \text{what evidence or derivation conflicts with it?}\\
\text{abduction} &:& \text{what hypothesis would explain the observation?}\\
\text{idk/unknown} &:& \text{what evidence is missing before answering?}
\end{array}
\label{eq:primitive_questions}
\end{equation}

This question view is not only interpretive. It determines the supervised labels
and losses attached to each head: primitive class, support status, entailment
state, role-to-primitive mapping, verifier likelihood, evidence pointer, IDK
action, tool intent, and schema mode. Causal interventions and register-level
probing are left as future interpretability work.

\subsection{Losses}
\label{subsec:losses}

The supervised objective is

\begin{equation}
\mathcal{L}
=
\lambda_{\mathrm{LM}}\LM
+ \sum_{h \in \mathcal{H}} \lambda_h
   \operatorname{CE}(y_h, f_h(X,E,P))
+ \lambda_{\mathrm{span}}\Lspan
+ \lambda_{\mathrm{iface}}\Liface
+ \beta \LKL.
\label{eq:total_loss}
\end{equation}

For ambiguous repeated logical atoms, strict first-occurrence span labels are
not semantically well-defined. Let $S^+$ be the set of all token positions that
are valid mentions of the same gold argument. The set-valued cross-entropy is

\begin{equation}
\mathcal{L}_{\mathrm{multi+}}(\ell, S^+)
=
\log \sum_j \exp(\ell_j)
-
\log \sum_{i \in S^+} \exp(\ell_i).
\label{eq:multipositive_ce}
\end{equation}

For joint span decoding, valid start/end pairs are restricted to a width band
$\mathcal{B}$:

\begin{equation}
(\hat{s},\hat{e})
=
\arg\max_{(s,e)\in\mathcal{B},\,e\ge s}
\left(\ell^{\mathrm{start}}_s + \ell^{\mathrm{end}}_e\right).
\label{eq:joint_span_decode}
\end{equation}

\subsection{Decoupled Interface Mode}
\label{subsec:interface_mode}

The usable checkpoint has two execution modes. Base mode uses the original
decoder blocks and measures the architecture:

\begin{equation}
y_{\mathrm{base}} = F_{\mathrm{D}}(x;\theta_{\mathrm{base}}).
\label{eq:base_mode}
\end{equation}

Interface mode swaps only the copied top-$K$ interface blocks:

\begin{equation}
y_{\mathrm{iface}} =
F_{\mathrm{D}}\!\left(x;\theta_{\mathrm{base}},
                     \theta_{\mathrm{iface}}^{(L-K+1:L)}\right).
\label{eq:interface_mode}
\end{equation}

The training invariant is

\begin{equation}
F_{\mathrm{D}}(x;\theta_{\mathrm{base}})
=
F_{\mathrm{D+iface}}(x;\theta_{\mathrm{base}})
\quad \text{in base mode},
\label{eq:base_identity}
\end{equation}

which is checked empirically as byte-identical primitive, role-to-primitive,
and LM metrics.
